\begin{abstract}
Coding agents can now write and debug substantial software, yet asking one to design a robot that
can be simulated mostly yields models that fail silently: parts float millimetres from where they
should attach, joint axes point the wrong way, links interpenetrate, and the masses are guesses.
These errors are invisible in the text of an MJCF or URDF file and surface only as odd behaviour
in simulation, long after they were made. We argue that the description language, not the agent,
is the bottleneck. We present \ms, a Python library in which a robot is one program and a part
can enter the robot only through a \emph{connection} to a part already in it; the library derives
every pose, joint anchor, axis and body frame from those connections and rejects a design with a
floating part at the line that caused it. Each build returns a machine-readable report of
verifiable properties: mass from the CAD solids, collisions at rest, the usable range of every
joint and what blocks it, and a physical test matched to the robot (standing on its actuators,
hovering, holding a payload). With this interface, six robots spanning an arm, a closed-loop
gripper, a quadrotor, a biped, a quadruped and a 29-DoF humanoid were rebuilt from public
references in 62--242 lines each, with every joint axis within 0.8\,mm of the reference's and
total mass within 6.2\% where the reference's mass is independent, and the resulting models train
locomotion policies directly.
\planned{one sentence with the headline number of the agent study (\cref{sec:agents}):
success rate of agents using \ms against agents writing MJCF/URDF or free-placement code CAD.}
\end{abstract}
